\cdcsection{Périmètre}
\label{sec:perimetre}

\cdcsubsection{Éléments inclus}

Pour garder un projet réalisable dans le cadre du module, le développement est organisé en deux versions successives. La version~1 constitue le socle livrable et utilisable. La version~2 l’enrichit sans le remettre en cause.

\cdcintertitre{Version 1 – Fonctionnalités essentielles}

\begin{cdcitems}
  \item Gestion des comptes, des boutiques et des réparateurs.
  \item Gestion et suivi des réparations.
  \item Photos et validation de l’état initial.
  \item Numéro de suivi client.
  \item Gestion des paiements en espèces.
  \item Enregistrement des paiements électroniques, si le prestataire est intégré.
  \item Calcul des commissions.
  \item Géolocalisation et recherche de boutiques.
\end{cdcitems}

\cdcintertitre{Version 2 – Fonctionnalités avancées}

\begin{cdcitems}
  \item Notifications WhatsApp automatisées.
  \item Diagnostic préliminaire intelligent.
  \item Comparaison avancée des boutiques et avis clients.
  \item Paiement électronique intégré avec confirmation automatique.
  \item Historique et gestion des garanties.
  \item Observatoire des prix et tableau de bord statistique avancé.
\end{cdcitems}

\cdcsubsection{Éléments exclus}

Les éléments suivants ne font pas partie du périmètre de Qribna.

\begin{cdcitems}
  \item Une application mobile native. Qribna est une plateforme web, et le client suit sa réparation sans installer d’application.
  \item La réalisation technique des réparations, qui reste sous la responsabilité de la boutique.
  \item Le diagnostic définitif. La plateforme fournit une estimation indicative, et seul le réparateur confirme la panne, le prix et le délai après examen de l’appareil.
  \item Le traitement direct des données bancaires. Les paiements électroniques passent exclusivement par un prestataire de paiement externe.
  \item La publication de données individuelles dans l’observatoire. Seules des statistiques agrégées sont diffusées.
\end{cdcitems}

\cdcsubsection{Utilisateurs concernés}

La plateforme s’adresse aux boutiques de réparation de téléphones du Maroc, en priorité aux petites boutiques, ainsi qu’à leurs clients. Quatre profils d’utilisateurs sont distingués.

\cdcintertitre{Le client}

Le client cherche une boutique, lui confie son téléphone et suit la réparation jusqu’au paiement.

\begin{cdcitems}
  \item Créer un compte ou rechercher une boutique.
  \item Décrire la panne de son téléphone.
  \item Consulter les prix estimés et les boutiques disponibles.
  \item Choisir une boutique et, éventuellement, réserver un rendez-vous.
  \item Suivre la réparation et recevoir des notifications.
  \item Choisir son moyen de paiement.
  \item Consulter ses reçus et son historique de réparations.
  \item Évaluer la boutique après la réparation.
\end{cdcitems}

\cdcintertitre{Le propriétaire de la boutique}

Le propriétaire administre sa boutique et pilote son activité.

\begin{cdcitems}
  \item Créer et gérer sa boutique.
  \item Ajouter et gérer les comptes des réparateurs.
  \item Consulter les réparations en cours et terminées.
  \item Définir les tarifs et les services proposés.
  \item Suivre les paiements, les revenus et les commissions.
  \item Consulter les statistiques et les performances de son activité.
\end{cdcitems}

\cdcintertitre{Le réparateur}

Le réparateur prend en charge les appareils et tient leur dossier à jour.

\begin{cdcitems}
  \item Enregistrer les téléphones reçus.
  \item Décrire les pannes et photographier l’état initial de l’appareil.
  \item Faire valider les informations par le client.
  \item Mettre à jour les étapes de la réparation.
  \item Enregistrer le coût final et les pièces utilisées.
  \item Indiquer le statut du paiement.
\end{cdcitems}

\cdcintertitre{L’administrateur de la plateforme}

L’administrateur veille à la qualité des informations publiées et au bon fonctionnement de l’ensemble.

\begin{cdcitems}
  \item Gérer les utilisateurs et les boutiques.
  \item Superviser les signalements et les avis.
  \item Contrôler les informations publiées.
  \item Assurer le bon fonctionnement général de la plateforme.
\end{cdcitems}
